% /root/monodromy, 2026-11-19. Root discovery/certificate; independent complete audit and rerun passed.
\subsection{An exact global contact in the degree-three norm family}
\label{subsec:ss-global-contact-witness}
\status{Proved and checked. The parent constructed and checked the exact
certificate. The monodromy agent independently audited its complete code,
recomputed the local data by truncated elliptic-curve addition, and reran
all four source-chart and all 144 collision checks, preserving the author's
outputs. This is one geometric witness, not a uniform full-group theorem.}

\begin{lemma}[A mixed contact with no other nonnodal collisions]
\label{lem:ss-global-contact03}
Let $E:y^2+y=x^3$, $X=E^2$, and
$L=\mathcal O_E(3O)\boxtimes\mathcal O_E(3O)$.
Over $\mathbf F_{64}=\mathbf F_2[z]/(z^6+z+1)$ encode
$\sum_{i=0}^5a_i z^i$ by the integer $\sum_{i=0}^5a_i2^i$.
Set $\omega=58$, so $\omega^2+\omega+1=0$, and let
\[
 i(x,y)=(x+1,y+x+\omega),\qquad b=1+i,\qquad
 \phi=\begin{pmatrix}1&b\\-b^\dagger&1\end{pmatrix}:X\to X.
\]
Then $\phi^\dagger\phi=3$ and $\ker\phi=\{(-bT,T):T\in E[3]\}$.
In the tensor basis $(1,x,y)\otimes(1,x,y)$, take
\[
 s=(38,14,10,47,47,43,30,4,8).
\]
The curve $C=(s=0)$ is geometrically smooth. Its norm
$Z=\phi(C)\in|L^3|$ has exactly one simple self-tacnode and otherwise
only ordinary nodes. The two preimages of the tacnode are
\[
 P=((17,52),(8,53)),\qquad Q=((18,35),(19,10)),
 \qquad Q-P=G=((58,59),(0,0)).
\]
The source Cartier kernel is the two-dimensional ambient kernel plus
one additional form whose values at these two translated tangent
vectors differ. Consequently the degree-three relative Hasse divisor
is smooth at $Z$.
\end{lemma}
\begin{proof}
The quaternion identity $i^\dagger=-i$, $i^2=-1$ gives
$b^\dagger b=2$, and hence the stated isogeny identities and kernel.
The exact norm line bundle is $L^3$, as proved in
Theorem~\ref{thm:ss-uniform-GL8-block}.
Write $S=(s_{ij})$ and $J=J_3\otimes J_3$, with $J_3$ reversing the
three basis vectors. The coefficient certificate gives
\[
 \det S=43\ne0,\qquad s^{[2]t}Js=0,\qquad\rank M(s)=8,
\]
and the literal identity
\[
 M(s)u=s^{[2]},\qquad
             u=(3,5,56,20,4,38,38,12,0).
\]
Here $M(s)$ is the root-removed Cartier multiplication matrix,
whose one-factor matrix is
\[
 M_E(a)=\begin{pmatrix}a_1&a_0&0\\a_2&0&a_0\\0&a_2&a_1\end{pmatrix}.
\]
Thus the induced Cartier map $W/ks\to W^{(2)}/ks^{(2)}$ has
one-dimensional kernel, represented by $u$.
Indeed $\ker M(s)=ks$ and $s^{[2]}$ lies in its image; quotienting
by that image line increases the kernel dimension by one before
quotienting the source by $ks$.

On the smooth source curve, the two invariant ambient one-forms
are Cartier-killed. Their adjunction boundary is the differential
of $\lambda_L=3\lambda_\Theta$, and is an isomorphism.
Therefore its full Cartier kernel has dimension three, and the
additional kernel form is represented by
$\operatorname{Res}(u\,dx_1\wedge dx_2/s)$, up to ambient correction.
Direct evaluation gives
\[
 s(P)=s(Q)=0,\qquad ds(P)=ds(Q)=(34,53),\qquad
                         u(P)+u(Q)=60\ne0.
\]
The derivatives use the invariant derivations
$\delta x=1$, $\delta y=x^2$ on each factor. The common tangent
vector $(53,34)$ makes the displayed difference precisely the
value difference of the additional differential; ambient corrections
have equal values on translated tangent vectors.
The independent truncated-series calculation, using the rational
elliptic addition law, finds an actual source-curve branch through
$P$ to order two and obtains
\[
                 s(\gamma(t))=O(t^3),\qquad
                 s(\gamma(t)+G)=55t^2+O(t^3).
\]
The nonzero coefficient proves simple contact, not higher tangency.

We describe the global exclusion, including its projective coverage.
The two charts $Z=1$ and $Y=1$ cover the cubic
$Y^2Z+YZ^2=X^3$. Both have equation $y^2+y=x^3$; the section
basis is $(1,x,y)$ in the first and $(y,x,1)$ in the second.
On all four product charts the ideal generated by the two elliptic
equations, $s$, and its two invariant derivatives has Groebner basis
$\{1\}$. Hence $C$ has no geometric singular point.

For every nonzero kernel translation $g$, an invertible projective
matrix for translation on each cubic was first computed. It was
then certified by exact rational addition identities modulo
$y^2+y+x^3$; agreement on the nine torsion points alone was not
used as a proof. For each of the eight $g$ and each of the four
product charts, the ideal
\[
       (F_1,F_2,s,t_g^*s,\det(ds,d(t_g^*s)))
\]
was checked. Its support is empty except at $P$ for $g=G$ and
at $Q$ for $g=-G$. For these two exceptions, reduction of
$(x_j-a_j)^{32}$ to zero in the Groebner basis proves that the
support over the algebraic closure is just the indicated point.
Frobenius is injective on an algebraically closed field, so this
power test does not merely exclude rational extra points.

Finally, for all twenty-eight unordered pairs $g\ne h$ of nonzero
kernel elements and on all four charts, the ideal
$(F_1,F_2,s,t_g^*s,t_h^*s)$ has Groebner basis $\{1\}$.
Thus no three distinct source points have the same image.
The eight translation matrices are genuine global translations,
so these $32+112=144$ checks exhaust all possible nonnodal
identifications. They also exclude invariance of $C$ by a nonzero
kernel translation. Since $C$ is a smooth connected ample divisor,
it is integral; $\phi|_C$ is therefore birational to its image.
Etaleness of $\phi$ gives no other singularities. The only nontransverse
identification is the certified simple contact; all other
identifications are transverse pairs. This proves the assertion
about the norm curve. Lemma~\ref{lem:ss-mixed-self-tacnode} now
proves smoothness of the relative Hasse divisor at its image.
\end{proof}

\begin{remark}[Retained discovery, failures and independent checks]
The parent's discovery and exact global verifiers are
\path{code/ss_contact_witness.py},
\path{code/ss_contact_witness_screened.py},
\path{code/verify_ss_contact_geometry.py}, and
\path{code/verify_ss_contact_collisions.py}; the successful input is
\path{outputs/ss_contact_witness_03.json}. The full successful outputs are
\path{outputs/verify_ss_contact_witness_03_geometry.json} and
\path{outputs/verify_ss_contact_witness_03_collisions.json}.
The third discovery used seed $202611191119$, and succeeded at its
first base pair and twenty-ninth section trial; its complete run
record is retained. The first two candidates failed the global test
because they had unwanted triple collisions. Their discovery and
failed verification outputs were retained; neither was promoted.

The independent local program
\path{agents/monodromy/code/verify_level3_contact03_local.py}
reconstructs the coefficient-field arithmetic, tensor matrix and
actual translated power series; it passed in
$0.0036840360007772688$ seconds.
The complete independent rerun of the four source charts passed in
$3.4523235300002852$ seconds and that of all $144$ collision ideals in
$72.10777428700021$ seconds. The parent's corresponding successful
runs took $3.2001894870008982$ and $74.52918819499973$ seconds.
The independent outputs are in the monodromy agent's output directory,
with the same verification basenames; the author's outputs were
preserved. The rerun driver and log are
\path{agents/monodromy/code/rerun_contact03_global_audit.py} and
\path{agents/monodromy/outputs/independent_contact03_global_run.log}.

This witness joins one pair-tangency locus to the source cubic $Q$.
It does not by itself establish irreducibility of the target Hasse
divisor: the two tangent-triple incidence orbits in the proposed
finite reduction still require their own global witnesses.
\end{remark}
